% Generated by experiments/scripts/export_latex.py -- do not edit by hand.
% Source: experiments/data/2-processed-tables-and-evaluations/revised-usd/robustness/robustness.json
% Source: experiments/data/0-usd-token-list-and-prices/constants.json
\begin{table}[htbp]
\centering
\caption[강건성 점검: 금액 기울기를 열 배 작게 하거나 크게 했을 때의 지표 가족 평균 켄달의 $\tau$.]{강건성 점검이에요. 금액 무게 $V$의 기울기 $b$를 열로 나누거나 열 배로 키웠을 때, 매칭 코호트(계약 $n=15{,}107$개)의 지표 가족 평균 켄달의 $\tau$를 주 분석 옆에 보여 줘요. 주 분석의 기울기는 금액이 중앙값일 때 $V=1/2$이 되게 정했어요. 그 중앙값은 송금이 \$200, 허락이 \$62예요. $b/10$은 그 점을 중앙값의 열 배로 옮기고, $10b$는 중앙값의 십분의 일로 옮겨요. 각 변형과, 같은 방법(엔도스랭크와 AWP)의 주 점수 사이의 순위 일치 $\tau$는 $b/10$에서 $0.973$과 $0.973$이고, $10b$에서 $0.968$과 $0.971$이에요. 변형에 대해서는 구간을 계산하지 않았어요.}
\label{tab:robustness-slope}
\fitwidth{%
\begin{tabular}{lccccccccc}
\toprule
방법 & \multicolumn{3}{c}{송금} & \multicolumn{3}{c}{허락} & \multicolumn{3}{c}{시빌} \\
\cmidrule(lr){2-4} \cmidrule(lr){5-7} \cmidrule(lr){8-10}
 & $b/10$ & $b$ & $10b$ & $b/10$ & $b$ & $10b$ & $b/10$ & $b$ & $10b$ \\
\midrule
EndorseRank & 0.157 & 0.158 & 0.159 & 0.588 & 0.585 & 0.580 & 0.006 & 0.007 & 0.009 \\
AWP & 0.577 & 0.572 & 0.567 & 0.193 & 0.193 & 0.191 & 0.397 & 0.395 & 0.393 \\
\bottomrule
\end{tabular}
}
\end{table}
